\chapter{Diseases of the Liver}
\label{ch:liver}
The liver is the largest internal organ, weighing approximately 1.5 kg, and performs over 500 vital functions including metabolism of carbohydrates, lipids, and proteins; synthesis of clotting factors and albumin; detoxification of drugs and toxins; bile production; and storage of vitamins and minerals. Hepatic diseases are diverse, ranging from steatosis and hepatitis to cirrhosis and hepatocellular carcinoma.

%-------------------------------------------------------------
\section{Alcoholic Liver Disease}
%-------------------------------------------------------------

%-------------------------------------------------------------

%-------------------------------------------------------------

%-------------------------------------------------------------

%-------------------------------------------------------------

%-------------------------------------------------------------

\label{sec:section:Alcoholic Liver Disease}
%-------------------------------------------------------------%-------------------------------------------------------------

\subsection{Alcoholic Liver Disease}
\label{sec:Alcoholic Liver Disease}
Alcoholic liver disease is a spectrum of hepatic injury due to excessive alcohol consumption that includes simple steatosis, alcoholic hepatitis, and cirrhosis. Simple steatosis is present in over 90\% of heavy drinkers. The condition results from direct hepatotoxicity of alcohol and its metabolites. Alcoholic hepatitis presents with jaundice, fever, hepatomegaly, and leukocytosis, carrying significant mortality (up to 40\% in severe cases); alcoholic cirrhosis develops in 10--20\% of heavy drinkers. Clinical features depend on the stage of disease. Diagnosis is based on history of significant alcohol intake (typically more than 50 g/day in men, more than 40 g/day in women for more than 5 years), clinical and laboratory findings, and exclusion of other causes. Management centers on complete alcohol abstinence, nutritional support, corticosteroids for severe alcoholic hepatitis, and liver transplantation for end-stage disease. Prognosis is favorable with abstinence but poor with continued drinking.

%-------------------------------------------------------------
\section{Cirrhosis}
%-------------------------------------------------------------

%-------------------------------------------------------------

Cirrhosis represents the irreversible end-stage result of chronic liver injury characterized by diffuse fibrosis, parenchymal necrosis, and regenerative nodule formation. Disruption of hepatic architecture increases portal vascular resistance, precipitating portal hypertension, splenomegaly, ascites, and variceal bleeding. Loss of synthetic and metabolic capacity leads to hepatic encephalopathy, coagulopathy, jaundice, and elevated risk of hepatocellular carcinoma.

%-------------------------------------------------------------

%-------------------------------------------------------------

%-------------------------------------------------------------

%-------------------------------------------------------------

\label{sec:section:Cirrhosis}
%-------------------------------------------------------------%-------------------------------------------------------------

\subsection{Ascites}
\label{sec:Ascites}
Ascites is the accumulation of fluid in the peritoneal cavity and is the most common complication of decompensated cirrhosis. The condition results from portal hypertension (sinusoidal pressure greater than 12 mmHg) and splanchnic arterial vasodilation, leading to neurohormonal activation (RAAS, sympathetic nervous system, ADH) and sodium and water retention. Clinical features include abdominal distension and discomfort. Diagnosis is by clinical examination (shifting dullness, fluid thrill) and abdominal ultrasound; diagnostic paracentesis shows a serum-ascites albumin gradient (SAAG) of 1.1 g/dL or greater, confirming portal hypertension. Management includes sodium restriction (less than 2 g/day), diuretics (spironolactone and furosemide), and large-volume paracentesis for tense ascites; refractory ascites may require transjugular intrahepatic portosystemic shunt (TIPS) or liver transplantation. Prognosis depends on the response to treatment and underlying liver disease severity; spontaneous bacterial peritonitis is a serious complication.

\subsection{Cirrhosis}
\label{sec:Cirrhosis}
Cirrhosis is the end stage of chronic liver disease characterized by diffuse fibrosis and formation of regenerative nodules that distort the hepatic architecture and vasculature, and is the 12th leading cause of death worldwide. The condition results from chronic liver injury due to viral hepatitis (B, C), alcohol, NAFLD/NASH, autoimmune hepatitis, primary biliary cholangitis, primary sclerosing cholangitis, or hereditary hemochromatosis. Cirrhosis may be compensated (no complications) or decompensated (development of ascites, variceal hemorrhage, hepatic encephalopathy, or jaundice). Portal hypertension drives many complications including esophageal and gastric varices, ascites, hepatorenal syndrome, and splenomegaly with hypersplenism. Clinical features depend on the stage and complications. Diagnosis is by liver biopsy, imaging (ultrasound, CT, MRI), and elastography (FibroScan); the MELD score predicts mortality and guides transplant listing. Management includes treatment of the underlying cause (antivirals, alcohol abstinence, weight loss), diuretics for ascites, beta-blockers or band ligation for variceal prophylaxis, lactulose and rifaximin for hepatic encephalopathy, and liver transplantation for end-stage disease. Prognosis depends on the stage and complications, with transplantation offering definitive treatment.

\subsection{Hepatic Encephalopathy}
\label{sec:Hepatic Encephalopathy}
Hepatic encephalopathy is a neuropsychiatric syndrome occurring in patients with advanced liver disease or portosystemic shunting, resulting from accumulation of neurotoxic substances (ammonia being the most important) that are normally cleared by the liver. The condition results from impaired hepatic clearance of neurotoxins, precipitated by infection, gastrointestinal bleeding, constipation, electrolyte imbalances, dehydration, and hepatorenal syndrome. Clinical features range from subtle cognitive impairment (grade 1) to confusion (grade 2), somnolence (grade 3), and coma (grade 4). Diagnosis is clinical, supported by elevated ammonia levels and exclusion of other causes of altered mental status. Management targets precipitating factors and reduces ammonia production and absorption with lactulose (titrated to 2--3 loose stools per day) and rifaximin. Prognosis is related to the severity of underlying liver disease.

%-------------------------------------------------------------
\section{Hepatic Neoplasms}
%-------------------------------------------------------------

%-------------------------------------------------------------

Hepatic neoplasms include primary benign lesions, primary malignant tumors, and secondary metastatic liver deposits originating from gastrointestinal or systemic malignancies. Primary liver cancers, dominated by hepatocellular carcinoma and cholangiocarcinoma, typically arise in the setting of chronic liver inflammation or cirrhosis. Multidisciplinary evaluation integrates multiphasic CT/MRI imaging, tumor biomarkers, surgical resection, local ablative techniques, and liver transplantation.

%-------------------------------------------------------------

%-------------------------------------------------------------

%-------------------------------------------------------------

%-------------------------------------------------------------

\label{sec:Hepatic Neoplasms}
%-------------------------------------------------------------%-------------------------------------------------------------

\subsection{Hepatocellular Carcinoma}
\label{sec:Hepatocellular Carcinoma}
Hepatocellular carcinoma is the most common primary liver cancer and the third leading cause of cancer-related death worldwide, arising predominantly in the setting of chronic liver disease and cirrhosis. Major risk factors include chronic HBV, chronic HCV, alcohol-related liver disease, NAFLD/NASH, and hereditary hemochromatosis. The condition results from malignant transformation of hepatocytes in the setting of chronic liver disease. Clinical features may include abdominal pain, weight loss, and hepatic decompensation. Diagnosis in cirrhotic patients can be made non-invasively with imaging (CT or MRI showing arterial enhancement with washout in the portal venous or delayed phase) using LI-RADS criteria; serum alpha-fetoprotein is a tumor marker. Staging uses the Barcelona Clinic Liver Cancer (BCLC) system. Management depends on stage: curative options include surgical resection, liver transplantation (Milan criteria), and radiofrequency ablation for early-stage disease; intermediate-stage disease is treated with transarterial chemoembolization (TACE); advanced disease is managed with systemic therapy (sorafenib, lenvatinib, atezolizumab plus bevacizumab). Prognosis depends on stage at diagnosis and treatment received.

%-------------------------------------------------------------
\section{Hepatitis}
%-------------------------------------------------------------

%-------------------------------------------------------------

Hepatitis refers to inflammation of the liver parenchyma caused by hepatotropic viruses, toxic substances, metabolic derangements, or autoimmune destruction. Acute hepatocellular injury causes hepatocellular necrosis, elevated serum aminotransferases, and impaired hepatic excretion. Chronic inflammation lasting beyond six months carries significant risk of progressive hepatic fibrosis, cirrhosis, and liver failure.

%-------------------------------------------------------------

%-------------------------------------------------------------

%-------------------------------------------------------------

%-------------------------------------------------------------

\label{sec:Hepatitis}
%-------------------------------------------------------------%-------------------------------------------------------------

\subsection{Autoimmune Hepatitis}
\label{sec:Autoimmune Hepatitis}
Autoimmune hepatitis is a chronic inflammatory liver disease characterized by autoimmune destruction of hepatocytes, predominantly affecting women. The condition results from autoimmune-mediated hepatocyte destruction. Type 1 (positive ANA or anti-smooth muscle antibodies) is more common and occurs at any age; type 2 (positive anti-LKM-1 antibodies) is rarer and typically affects children and young women. Clinical features range from asymptomatic elevation of aminotransferases to acute hepatitis or fulminant failure. Diagnosis requires elevated transaminases, hypergammaglobulinemia (IgG), characteristic autoantibodies, and liver biopsy showing interface hepatitis with plasma cell infiltration. Management involves corticosteroids (prednisone) combined with azathioprine to induce and maintain remission. Prognosis is good with treatment, although lifelong immunosuppression is often necessary.

\subsection{Hepatitis A}
\label{sec:Hepatitis A}
Hepatitis A is an acute, self-limited viral hepatitis transmitted via the fecal-oral route through contaminated food or water or close contact with infected individuals. The condition results from infection with hepatitis A virus (HAV), with an incubation period of 15--50 days. Most infections in children are asymptomatic; adults present with malaise, anorexia, nausea, vomiting, right upper quadrant pain, jaundice, and dark urine. Clinical features include malaise, anorexia, nausea, vomiting, right upper quadrant pain, jaundice, and dark urine. Diagnosis is based on elevated serum aminotransferases (ALT, AST) and confirmed by anti-HAV IgM antibodies. Management is supportive, as no specific antiviral treatment exists; prevention is through vaccination (inactivated virus) and hygiene measures. Prognosis is excellent, as HAV does not cause chronic infection or cirrhosis; fulminant hepatic failure is rare (0.1--0.3\%).

\subsection{Hepatitis B}
\label{sec:Hepatitis B}
Hepatitis B is a major global infection caused by hepatitis B virus (HBV) that affects approximately 296 million chronically infected individuals worldwide. Transmission occurs through percutaneous exposure (needle sharing, needlestick injuries), sexual contact, and perinatal transmission. The condition results from HBV infection, with chronic infection defined by persistence of HBsAg for more than 6 months; chronic hepatitis B carries risks of cirrhosis (15--40\%) and hepatocellular carcinoma (2--5\% per year in cirrhotics). Clinical features of acute infection range from subclinical to fulminant hepatic failure; the natural history follows a spectrum of immune tolerance, immune active phase, inactive carrier state, and reactivation. Diagnosis involves serological testing (HBsAg, anti-HBs, HBeAg, anti-HBe, anti-HBc IgM/IgG) and HBV DNA quantification. Management with nucleos(t)ide analogues (entecavir, tenofovir) or pegylated interferon-alpha is indicated for active disease. Prognosis is improved with antiviral therapy; vaccination is highly effective in preventing infection.

\subsection{Hepatitis C}
\label{sec:Hepatitis C}
Hepatitis C is a viral infection caused by hepatitis C virus (HCV) that affects approximately 58 million people globally and is a leading cause of cirrhosis, liver failure, and hepatocellular carcinoma. Transmission is primarily through blood exposure (intravenous drug use, blood transfusions before screening, needlestick injuries). The condition results from HCV infection, with approximately 75--85\% of acute cases progressing to chronic infection. Chronic hepatitis C follows an indolent course, with 10--20\% developing cirrhosis over 20--30 years; extrahepatic manifestations include cryoglobulinemia, membranoproliferative glomerulonephritis, and B-cell lymphoma. Clinical features of acute infection are usually asymptomatic. Diagnosis involves anti-HCV antibody testing followed by HCV RNA quantification and genotyping. Management with direct-acting antivirals (DAAs) achieves sustained virological response rates exceeding 95\% with 8--12 weeks of oral therapy. Prognosis is excellent with treatment; the WHO has set a goal of hepatitis C elimination by 2030.

\subsection{Hepatitis D}
\label{sec:Hepatitis D}
Hepatitis D is an infection caused by hepatitis D virus (HDV), a defective RNA virus that requires co-infection or superinfection with HBV (using HBsAg for its envelope), affecting approximately 5\% of HBV-infected individuals. The condition results from HDV infection in the presence of HBV. Co-infection results in more severe acute hepatitis and a higher risk of fulminant hepatic failure than HBV alone; superinfection in chronic HBV carriers accelerates progression to cirrhosis and HCC. Clinical features include more severe hepatitis compared with HBV monoinfection. Diagnosis is by anti-HDV antibody testing and HDV RNA quantification. Management involves pegylated interferon-alpha with limited efficacy; bulevirtide has been approved in Europe as an entry inhibitor. Prognosis is worse than HBV monoinfection, although prevention of HBV through vaccination prevents HDV infection.

\subsection{Hepatitis E}
\label{sec:Hepatitis E}
Hepatitis E is an inflammatory liver disease caused by hepatitis E virus (HEV, Orthohepevirus A), a non-enveloped single-stranded RNA virus in the Hepeviridae family, with four major genotypes: genotypes 1 and 2 cause waterborne epidemics in developing countries, and genotypes 3 and 4 cause zoonotic infection (primarily from swine) in developed countries. An estimated 20 million HEV infections occur annually, with 3.3 million symptomatic cases and 44,000 deaths. The condition results from fecal-oral transmission (genotypes 1, 2 through contaminated water) or consumption of undercooked pork or game meat (genotypes 3, 4). HEV replicates in hepatocytes, causing hepatocellular injury and necrosis. Clinical features of acute hepatitis E include jaundice, malaise, anorexia, abdominal pain, hepatomegaly, fever, and elevated transaminases, and are usually self-limited over 4--6 weeks. A distinctive feature is severe disease in pregnancy (genotypes 1, 2), with mortality of 20--30\% in the third trimester, high rates of fulminant hepatic failure, and adverse fetal outcomes (preterm delivery, stillbirth). Chronic HEV infection occurs in immunocompromised individuals (organ transplant recipients, HIV, hematologic malignancy) with genotypes 3, 4, leading to rapidly progressive cirrhosis. Diagnosis is based on serum anti-HEV IgM and HEV RNA detection (PCR). Management is supportive for acute hepatitis E; ribavirin for chronic HEV (12 weeks); reduction of immunosuppression is first-line for transplant recipients. A recombinant HEV vaccine (Hecolin) is licensed in China. Prognosis is excellent for acute HEV in immunocompetent hosts; chronic HEV is treatable with ribavirin.

%-------------------------------------------------------------
\section{Hereditary and Metabolic Liver Diseases}
%-------------------------------------------------------------

%-------------------------------------------------------------

%-------------------------------------------------------------

%-------------------------------------------------------------

%-------------------------------------------------------------

%-------------------------------------------------------------

\label{sec:Hereditary and Metabolic Liver Diseases}
%-------------------------------------------------------------%-------------------------------------------------------------

\subsection{Hereditary Hemochromatosis}
\label{sec:Hereditary Hemochromatosis}
Hereditary hemochromatosis is a group of genetic disorders characterized by excessive intestinal iron absorption and progressive iron deposition in tissues, leading to organ damage; the most common form is type 1 (HFE-related HH), an autosomal recessive disorder caused by mutations in the HFE gene (predominantly C282Y homozygosity), affecting approximately 1 in 200--300 individuals of Northern European descent, though clinical disease develops in only 10--30\% of C282Y homozygotes. The condition results from excessive iron absorption leading to iron deposition in the liver, heart, pancreas, joints, skin, and endocrine glands. Clinical manifestations include hepatomegaly, cirrhosis, hepatocellular carcinoma, diabetes mellitus, cardiomyopathy, arthropathy, hypogonadism, and skin hyperpigmentation. Diagnosis involves elevated serum transferrin saturation (greater than 45\%), elevated serum ferritin, and HFE genotyping; liver biopsy or MRI T2* quantification assesses iron burden. Management involves regular phlebotomy, which is highly effective in preventing complications if initiated early. Prognosis is excellent with early diagnosis and treatment but poor if cirrhosis or cardiomyopathy has developed; first-degree relatives should be screened.

%-------------------------------------------------------------
\section{Metabolic and Fatty Liver Disease}
%-------------------------------------------------------------

%-------------------------------------------------------------

%-------------------------------------------------------------

%-------------------------------------------------------------

%-------------------------------------------------------------

%-------------------------------------------------------------

\label{sec:Metabolic and Fatty Liver Disease}
%-------------------------------------------------------------%-------------------------------------------------------------

\subsection{Non-Alcoholic Fatty Liver Disease}
\label{sec:Non-Alcoholic Fatty Liver Disease}
Non-alcoholic fatty liver disease (NAFLD), now increasingly termed metabolic dysfunction-associated steatotic liver disease (MASLD), is the most common chronic liver disease worldwide, affecting approximately 25\% of the global population, and is strongly associated with metabolic syndrome, obesity, type 2 diabetes mellitus, and dyslipidemia. The condition results from hepatic steatosis due to insulin resistance and metabolic dysfunction. The spectrum ranges from simple steatosis to non-alcoholic steatohepatitis (NASH/MASH), which involves hepatocyte injury, inflammation, and fibrosis, and can progress to cirrhosis and hepatocellular carcinoma. Clinical features are often absent, with elevated transaminases detected incidentally. Diagnosis is by imaging (ultrasound, CT, MRI-PDFF) showing hepatic steatosis and exclusion of other causes; liver biopsy remains the gold standard for staging. Management involves weight loss through diet and exercise (7--10\% body weight reduction), management of metabolic risk factors, vitamin E or pioglitazone for non-diabetic NASH, and resmetirom for NASH with fibrosis. Prognosis depends on the stage of fibrosis.

\subsection{Wilson's Disease}
\label{sec:Wilson's Disease}
Wilson's disease is an autosomal recessive disorder of copper metabolism caused by mutations in the ATP7B gene, with a prevalence of approximately 1 in 30,000. The condition results from impaired biliary copper excretion leading to accumulation of copper in the liver, brain, cornea, and other organs. Hepatic manifestations range from asymptomatic transaminase elevation to acute hepatitis, chronic hepatitis, or cirrhosis; neuropsychiatric features include tremor, dysarthria, dystonia, parkinsonism, personality changes, and depression. Kayser-Fleischer rings (copper deposition in Descemet's membrane of the cornea) are a hallmark finding. Diagnosis involves low serum ceruloplasmin (less than 20 mg/dL), elevated 24-hour urinary copper excretion, and hepatic copper quantification on liver biopsy. Management includes copper chelation (penicillamine, trientine), zinc therapy to block copper absorption, and liver transplantation in fulminant or decompensated disease; screening of first-degree relatives is essential. Prognosis is excellent with early diagnosis and lifelong treatment.

%-------------------------------------------------------------
\section{Vascular Liver Diseases}
%-------------------------------------------------------------

%-------------------------------------------------------------

%-------------------------------------------------------------

%-------------------------------------------------------------

%-------------------------------------------------------------

%-------------------------------------------------------------

\label{sec:Vascular Liver Diseases}
%-------------------------------------------------------------%-------------------------------------------------------------

\subsection{Budd-Chiari Syndrome}
\label{sec:Budd-Chiari Syndrome}
Budd-Chiari syndrome is a hepatic venous outflow obstruction at any level from the small hepatic veins to the junction of the inferior vena cava and right atrium. The most common cause is myeloproliferative neoplasms (particularly polycythemia vera); other causes include antiphospholipid syndrome, paroxysmal nocturnal hemoglobinuria, and oral contraceptive use. The condition results from obstruction of hepatic venous outflow, leading to hepatic congestion, sinusoidal hypertension, and hepatocyte ischemia. Clinical features range from asymptomatic to fulminant hepatic failure; classic features include acute abdominal pain, hepatomegaly, and ascites. Diagnosis is by Doppler ultrasound, CT, or MR venography showing absent or abnormal hepatic venous flow. Management includes anticoagulation, catheter-directed thrombolysis, TIPS for decompression, and liver transplantation for progressive liver failure. Prognosis depends on the cause and severity.

\subsection{Portal Hypertension}
\label{sec:Portal Hypertension}
Portal hypertension is defined as a hepatic venous pressure gradient (HVPG) of 5 mmHg or greater, with clinically significant portal hypertension at HVPG of 10 mmHg or greater; the most common cause is cirrhosis. The condition results from increased resistance to portal blood flow. Consequences include esophageal and gastric varices, ascites, portosystemic collaterals, splenomegaly, and hepatorenal syndrome. Clinical features depend on the complications present. Diagnosis is confirmed by HVPG measurement via right heart catheterization; non-invasive assessment includes platelet count, transient elastography, and endoscopy. Management includes non-selective beta-blockers (propranolol, nadolol) and endoscopic variceal band ligation for primary and secondary prophylaxis of variceal hemorrhage. Prognosis depends on the underlying liver disease and response to treatment.

\subsection{Primary Biliary Cholangitis}
\label{sec:Primary Biliary Cholangitis}
Primary biliary cholangitis (PBC, formerly primary biliary cirrhosis) is a chronic, progressive autoimmune liver disease characterized by granulomatous destruction of intrahepatic interlobular bile ducts, occurring predominantly in women aged 40--60. The condition results from autoimmune-mediated destruction of biliary epithelial cells. Clinical features include fatigue, severe generalized pruritus (often preceding jaundice), hyperpigmentation, and xanthomas. Diagnosis is confirmed by elevated serum alkaline phosphatase, positive antimitochondrial antibodies (AMA, present in $>95\%$ of patients), and liver biopsy showing nonsuppurative destructive cholangitis. Management with ursodeoxycholic acid (UDCA) slows disease progression; obeticholic acid is used in UDCA non-responders; liver transplantation is definitive for end-stage disease. Prognosis is significantly improved with UDCA treatment.

\subsection{Primary Sclerosing Cholangitis}
\label{sec:Primary Sclerosing Cholangitis}
Primary sclerosing cholangitis (PSC) is a chronic cholestatic liver disease characterized by inflammation, fibrotic stricturing, and obliteration of intrahepatic and extrahepatic bile ducts, strongly associated with ulcerative colitis (present in 70--80\% of PSC patients). The condition results from immune-mediated bile duct injury leading to concentric onion-skin fibrosis around bile ducts. Clinical features include fatigue, pruritus, right upper quadrant pain, recurrent acute cholangitis, and progressive jaundice. Diagnosis is established by magnetic resonance cholangiopancreatography (MRCP) or ERCP showing characteristic beaded appearance of bile ducts. Management is supportive; endoscopic balloon dilation or stenting is used for dominant strictures; liver transplantation is the only curative option. PSC carries a 10--15\% lifetime risk of cholangiocarcinoma. Prognosis is guarded without liver transplantation.

